%=======================================================================
%  Regime-Aware Cooperation in Multi-Echelon Inventory Control:
%  Curriculum Pretraining and a Zero-Shot LLM Supervisor
%  IEEE Conference format (IEEEtran) -- 6 page target
%=======================================================================
\documentclass[conference]{IEEEtran}
\IEEEoverridecommandlockouts

\usepackage{cite}
\usepackage{amsmath,amssymb,amsfonts}
\usepackage{graphicx}
\usepackage{textcomp}
\usepackage{booktabs}
\usepackage{multirow}
\usepackage{array}
\usepackage{xurl}

% Tighten float, caption and bibliography spacing to fit the page limit.
\setlength{\textfloatsep}{7pt plus 2pt minus 2pt}
\setlength{\floatsep}{7pt plus 2pt minus 2pt}
\setlength{\intextsep}{7pt plus 2pt minus 2pt}
\setlength{\dblfloatsep}{7pt plus 2pt minus 2pt}
\setlength{\dbltextfloatsep}{7pt plus 2pt minus 2pt}
\setlength{\abovecaptionskip}{4pt}
\setlength{\belowcaptionskip}{0pt}
\def\IEEEbibitemsep{0pt plus 0.5pt}

\begin{document}

\title{Regime-Aware Cooperation in Multi-Echelon Inventory Control:\\
Curriculum Pretraining and a Zero-Shot LLM Supervisor}

\author{
\IEEEauthorblockN{Chinmayee S Bharadwaj\IEEEauthorrefmark{1},
Taran Thimmaiah\IEEEauthorrefmark{1},
Alvin S Thottappilly\IEEEauthorrefmark{1},
Manoj Kumar\IEEEauthorrefmark{1}}
\IEEEauthorblockA{\IEEEauthorrefmark{1}Department of Computer Science and Engineering (Data Science)\\
B.M.S. College of Engineering, Bengaluru 560085, India\\
Email: \{chinmayee.cd23, taranthimmaiah.cd23, alvins.cd23, manojkumar.cds\}@bmsce.ac.in}
}

\maketitle

%=======================================================================
\begin{abstract}
Multi-agent deep reinforcement learning produces competitive decentralized
ordering policies for multi-echelon inventory systems. A central design
choice is the cooperation weight $\alpha$, which blends an actor's local
cost with the system cost in its reward; prior work tunes it offline and
freezes it. We instead train one HAPPO policy across a curriculum of three
cooperation regimes, $\alpha \in \{0.100, 0.375, 0.650\}$, sampled
$70/15/15$ per episode, and let a supervisor reassign a per-echelon $\alpha$
vector every $N=20$ periods. Evaluating on 280 frozen traces spanning
in-distribution demand and three disruption regimes, benchmarked against a
reproduction of the authors' released implementation, we find the curriculum
yields
robustness that is regime-dependent rather than uniform, and yields it at
execution time without run-time switching. Under a pulse shock it cuts
per-period cost by $12.8\%$ and under a demand collapse by $6.9\%$, both
with $\alpha$ held fixed at execution time; in distribution it matches the
specialist on cost while raising customer fill rate by $7.1$ points. Under
a sustained step-up it loses: with orders capped at $20$ units the chain is
capacity-bound, the buffer the curriculum teaches cannot clear the
downstream bottleneck, and cost rises $6.8\%$. Ablating the supervisor on
identical frozen weights, a language model reading window telemetry matches
a hand-tuned rule base to within $3$ cost units at $n=100$: expert heuristic
quality without manual threshold tuning and an auditable rationale
per decision. Switching itself pays most under contraction.
We report a fixed-$\alpha$ execution arm throughout, which makes the two
contributions separable, and recommend it in comparable studies.
\end{abstract}

\begin{IEEEkeywords}
multi-echelon inventory management, multi-agent reinforcement learning,
large language models, curriculum learning, zero-shot generalization,
bullwhip effect, explainability
\end{IEEEkeywords}

%=======================================================================
\section{Introduction}

Inventory decisions in a serial supply chain are coupled: a retailer's order
is demand for a warehouse, whose order is demand for a distributor. Because
each tier sees only its neighbour, small perturbations amplify upstream.
This is the bullwhip effect~\cite{lee1997}, which Sterman's beer-game
experiments showed arises from feedback misperception even under stationary
demand~\cite{sterman1989}. Classical multi-echelon theory gives elegant
results under i.i.d.\ demand and simple
topologies~\cite{clark1960,shang2003,zipkin2000}, assumptions rarely met in
practice.

Deep reinforcement learning is an established
alternative~\cite{boute2022}, though much of it assumes a central
controller with full visibility~\cite{gijsbrechts2021,harsha2022}, which
conflicts with the information asymmetry between real supply chain
partners~\cite{vosooghidizaji2020}. Liu \emph{et al.}~\cite{liu2022}
instead formulate the chain as a partially observable Markov game and
apply heterogeneous-agent PPO (HAPPO)~\cite{kuba2021} under centralized
training with decentralized execution. Their notable finding concerns the
reward: each actor minimizes a convex combination
\begin{equation}
r_t^i = -\big[\alpha\, C_t^i + (1-\alpha)\, C_t^{\text{total}}\big]
\label{eq:basereward}
\end{equation}
of its own cost and the system cost, and performance peaks neither at
$\alpha=0$ nor $\alpha=1$ but at an interior point, $0.375$ for a
four-echelon chain. Low $\alpha$ makes the policy-gradient target unbiased
but noisy, since the reward absorbs every other actor's stochasticity;
high $\alpha$ makes it clean but misaimed.

Our starting observation is that this balance is a property of the operating
regime, not of the chain. Under calm demand, variance dominates and a
moderately self-interested $\alpha$ is right; under a shock, suppressing
amplification calls for a smaller one; when a tier is deeply backlogged,
diluting accountability across four tiers delays the response and calls for
a larger one. A scalar fixed before deployment cannot express all three.

We therefore make $\alpha$ a control input. A single policy is trained
across three regimes under a skewed curriculum, and a supervisor reassigns a
per-echelon $\alpha$ vector every $20$ periods from aggregated telemetry.
The contributions are:

\begin{enumerate}
\item \textbf{A cooperation-weight curriculum buys regime-dependent
robustness.} Training one policy across three cooperation regimes cuts cost
by $12.8\%$ under a pulse shock and $6.9\%$ under a demand collapse relative
to a reproduced single-$\alpha$ specialist, and raises in-distribution fill
rate by $7.1$ points at cost parity. All of this holds with $\alpha$ fixed
at execution time, so the gain requires no run-time machinery. Under a
sustained step-up that drives demand to the order cap it instead costs
$6.8\%$ more, and we identify the mechanism.

\item \textbf{An LLM supervisor reproduces hand-tuned heuristics zero-shot.}
At $n=100$ the paired difference between the language model and a rule base
tuned against this environment is $-0.37$ with interval $[-3.07, +2.33]$:
an equivalence result bounding any true gap to roughly $\pm 3$ cost units.
The model is given the topology and the regimes in a structured telemetry prompt
without manual policy threshold rules, and logs a rationale for every decision.

\item \textbf{A fixed-$\alpha$ arm separates the two.} Executing the trained
policy at a constant $\alpha$ isolates what training contributes from what
switching contributes. Switching pays most under contraction, where rules
beat a static $\alpha$ by $4.77$; reporting the fixed arm is what makes that
attribution possible.
\end{enumerate}

%=======================================================================
\section{Related Work}

\textbf{DRL for inventory control.} Gijsbrechts
\emph{et al.}~\cite{gijsbrechts2021} apply A3C to lost-sales and
multi-echelon problems; Vanvuchelen \emph{et al.}~\cite{vanvuchelen2020}
use PPO~\cite{schulman2017} for joint replenishment; Harsha
\emph{et al.}~\cite{harsha2022} combine mathematical programming with
policy iteration. All use one learner prescribing actions for every
node; Oroojlooyjadid \emph{et al.}~\cite{oroojlooy2022} decentralize, but
for one actor only. Boute \emph{et al.}~\cite{boute2022} identify
generalization beyond the training distribution as the open problem, which
is our target.

\textbf{Multi-agent RL.} Centralized training with decentralized execution
resolves the non-stationarity of independent learners by conditioning the
critic on joint information while keeping actors local~\cite{lowe2017}.
MAPPO~\cite{yu2021} shows PPO with a centralized value function is a
strong cooperative baseline; HAPPO~\cite{kuba2021} adds a sequential
update scheme restoring monotonic improvement for heterogeneous agents.
We retain HAPPO unchanged so our comparison isolates the control
architecture rather than the learner.

\textbf{Curriculum and generalization.} Training on a distribution of tasks
is a standard route to policies that survive unseen conditions: curriculum
learning~\cite{bengio2009} orders that distribution and domain
randomization~\cite{tobin2017} widens it, trading in-distribution
optimality for transfer. Kirk \emph{et al.}~\cite{kirk2023} note that most
zero-shot gains come from varying the training distribution rather than the
architecture. Ours varies neither dynamics nor task but the \emph{reward
parameterization}, and the trade proves displaced rather than absent:
robustness under transient and contractionary disruption comes free in
distribution, and is paid for under sustained expansion.

\textbf{Language models as controllers.} LLMs have been used to sequence
grounded robot skills~\cite{ahn2022}, to interleave reasoning with
acting~\cite{yao2023}, as proxy reward functions~\cite{kwon2023}, to
author reward code~\cite{ma2024}, and as black-box
optimizers driven by natural-language problem
descriptions~\cite{yang2024}. Those works modify the reward before or during
training; we do neither. The reward is fixed, the policy frozen, and the
model adjusts only the scalar $\alpha$ the policy already takes as input.
Our supervisor is one arm of an ablation rather than the proposed
mechanism, which lets us report what the language model does and does not
contribute.

%=======================================================================
\section{Problem Formulation}

\subsection{Environment}
We model a four-echelon serial chain, Manufacturer $\rightarrow$
Distributor $\rightarrow$ Warehouse $\rightarrow$ Retailer $\rightarrow$
Customer, indexed $i \in \{0,\dots,3\}$ from the retailer upstream, as a
decentralized partially observable Markov game~\cite{littman1994}. Each
tier has lead time $L=4$, so a production decision needs $16$ periods to
reach customer-facing inventory; this delay governs several design
choices below.

At period $t$, tier $i$ faces effective demand $D_t^0 = D_t + B_{t-1}^0$
for the retailer and $D_t^i = a_t^{i-1} + B_{t-1}^i$ upstream, where
$a_t^{i-1}$ is the downstream order and $B$ is carried backlog. With
deficit $U_t^i = D_t^i - (I_{t-1}^i + P_{t-1}^{i,1})$, inventory and
backlog update as $I_t^i = \max(0,-U_t^i)$, $B_t^i = \max(0,U_t^i)$, and
the pipeline advances one slot per period. Actions are discrete,
$a_t^i \in \{0,\dots,20\}$. Costs are symmetric, $h=b=1$, so
$C_t^i = I_t^i + B_t^i$.

\subsection{Reward and the cooperation weight}
The published reward~\eqref{eq:basereward} combines an actor's own cost
with the \emph{total} system cost. The authors' released
implementation~\cite{liu2022code} instead combines it with the
\emph{mean}: in code the system term is \texttt{np.mean(reward)} where
\eqref{eq:basereward} is a sum. Our reproduced baseline and our own policy
both follow the released implementation, so throughout this paper
\begin{equation}
r_t^i = -\Big[\alpha_t^i\, C_t^i + (1-\alpha_t^i)\,\tfrac{1}{M} C_t^{\text{total}}\Big],
\quad M = 4 .
\label{eq:reward}
\end{equation}
The two forms are not equivalent: expanding
$C_t^{\text{total}} = C_t^i + \sum_{j \neq i} C_t^j$, the published reward
weights own cost by $1$ and each other actor by $(1-\alpha)$, whereas
\eqref{eq:reward} weights them by $\alpha + (1-\alpha)/M$ and
$(1-\alpha)/M$. Since the repository is designated by~\cite{liu2022} as its
official implementation and applies the mean unconditionally, the reported
$\alpha^{\ast}=0.375$ is expressed in that parameterization. The
bias--variance trade-off holds under either form, the ratio of other-cost to
own-cost weight decreasing monotonically in $\alpha$.

The modification we introduce is that $\alpha_t^i$ carries a time and an
agent index. It is a control input, not a hyperparameter.

\subsection{Observations}
The decentralized actor observation is
$\mathbf{o}_t^i = [\,I_t^i, B_t^i, D_t^{\text{in},i}, P_t^{i,1..4},
\alpha_t^i\,] \in \mathbb{R}^8$, each component normalized by the order
cap. The reference implementation omits the last component and uses
$\mathbb{R}^7$, its $\alpha$ being constant. The centralized critic receives the
$28$-dimensional concatenation of all four tiers' physical features and
never observes $\alpha$: the value function should estimate physical
cost-to-go, and a control variable that changes mid-episode would make its
target non-stationary for reasons unrelated to the inventory dynamics.

%=======================================================================
\section{Regime-Aware Cooperation}

\subsection{Three regimes}
We instantiate three regimes, chosen for what each does
to~\eqref{eq:reward} rather than by search.
\textsc{Cruiser} ($\alpha = 0.375$) is the value~\cite{liu2022} reports as
optimal for a four-echelon chain and the operating point of our reproduced
baseline; it is the anchor the controller must not give up.
\textsc{Shock Absorber} ($\alpha = 0.100$) pushes $90\%$ of the objective
onto system cost, so tiers internalize the downstream consequences of
their own order variance. This is the setting for damping amplification and
draining stock accumulated after a shock.
\textsc{Backlog Clearer} ($\alpha = 0.650$) inverts the weighting,
sharpening the gradient an individual tier receives from its own stockout,
which is what a starved tier needs when the others are not the bottleneck.

Discrete points are used rather than a continuous $\alpha$ because a
language model is a reliable categorical reasoner and an unreliable
continuous one, and because concentrating the curriculum on few, widely
separated modes preserves capacity at the operating point the chain occupies
most of the time. Three modes over four tiers still span $3^4 = 81$
configurations.

\subsection{Skewed curriculum}
We train all three operational regimes into a single network policy. At the
start of each training episode the environment randomly samples an operating
mode with probability $70\%$ for \textsc{Cruiser} ($\alpha=0.375$), $15\%$
for \textsc{Shock Absorber} ($\alpha=0.100$) and $15\%$ for \textsc{Backlog
Clearer} ($\alpha=0.650$). The selected $\alpha$ is appended directly to each
agent's observation vector as an additional input feature, making this a
curriculum in the sense of~\cite{bengio2009} over reward parameterizations
rather than over task difficulty.

Rather than training the modes in separate sequential phases, we interleave
them randomly across episodes. This prevents catastrophic forgetting,
ensuring the network does not lose its normal steady-state skills when
exposed to emergency modes. The $70/15/15$ split is deliberate: in real
operations a supply chain spends most of its time in calm conditions and
enters emergency states only occasionally, so a balanced $33/33/33$ split
would spend most gradient updates on regimes the chain rarely occupies. We
did not ablate the split, so the ratio is a design choice rather than a
tuned optimum.

\subsection{Dual-timescale execution}
Fig.~\ref{fig:arch} shows the system. The inner loop is standard
decentralized execution: every period, each actor reads its local
observation and emits an order, with no communication. The outer loop
activates at $t = 20, 40, \dots$: the environment aggregates the preceding
window into demand mean and spread, per-echelon backlog, on-hand inventory,
window cost and bullwhip amplification, and a supervisor
returns a mode with a per-echelon vector
$[\alpha_R, \alpha_W, \alpha_D, \alpha_M]$, written into the actors'
observations for the next window. The cadence is set by physics: a
replenishment wave takes $16$ periods to cross the chain, so intervening
more often would judge a decision before its effects arrive.

\begin{figure}[t]
\centering
\includegraphics[width=\columnwidth]{figures/architecture.pdf}
\caption{System architecture. Phase 1 trains one HAPPO policy across three
cooperation regimes under a $70/15/15$ curriculum; the centralized critic
sees the joint physical state but never $\alpha$. The weights are then
frozen. Phase 2 runs decentralized actors every period (solid) while a
supervisor (static, rule-based, or an LLM) reassigns a per-echelon
$\alpha$ vector every $20$ periods (dashed).}
\label{fig:arch}
\end{figure}

\subsection{Three supervisors}
Because the same frozen weights can be driven by any policy over $\alpha$,
the supervisor is an ablation dimension rather than a fixed component. We
evaluate three, all reading the same telemetry at the same cadence.
\textbf{Static} holds $\alpha = 0.375$ at every tier for the whole
episode, isolating what the curriculum alone contributes.
\textbf{Rule-based} applies the hand-tuned priority ladder and thresholds in
Table~\ref{tab:rules}, developed against this environment.
\textbf{LLM} receives window telemetry formatted as structured text along with
status indicators for each tier, and queries a frontier language model, the preview
endpoint \texttt{gpt-5.6-luna}. The model returns a structured JSON reply giving a
target mode, the per-echelon $\alpha$ vector, a diagnosis and an operational rationale. Each assigned
$\alpha$ is snapped to the nearest discrete anchor in $\{0.100, 0.375, 0.650\}$. Failed
calls or schema errors are retried up to $5$ times before defaulting to the rule-based
fallback (zero fallbacks occurred across all $280$ benchmark traces). The prompt states that
raising $\alpha$ during a spike produces no immediate inventory, only upstream order variance;
without it the model reaches for aggressive settings whenever demand rises.

\begin{table}[t]
\caption{Rule-based supervisor priority ladder. Evaluated in top-to-bottom order every $N=20$ periods. The assigned column lists $\alpha$ for retailer / warehouse / distributor / manufacturer; ``all'' means one value across every echelon.}
\label{tab:rules}
\centering
\footnotesize
\setlength{\tabcolsep}{4pt}
\begin{tabular}{@{}>{\raggedright\arraybackslash}p{2.6cm}>{\raggedright\arraybackslash}p{1.9cm}>{\raggedright\arraybackslash}p{3.5cm}@{}}
\toprule
Trigger Condition & Assigned $\alpha$ & Operational Intent \\
\midrule
1. Mean demand $\ge 14$ \textbf{or} $B_R{+}B_W > 40$ (max 2 consecutive windows) & $.65/.65/.10/.65$ & Multi-tier recovery: pull downstream, saturate factory, distributor frictionless \\
\addlinespace[3pt]
2. $\ge 2$ demand spikes, std $> 3.5$, \textbf{or} Bullwhip Amp $> 120\%$ & all $.10$ & Shock absorber: damp ripple and upstream variance surge \\
\addlinespace[3pt]
3. Total backlog $> 4$ \textbf{or} mean backlog $> 2$ (non-consecutive) & $.65/.10/.10/.10$ & Targeted hybrid: retailer pulls, upstream flushes stock without hoarding \\
\addlinespace[3pt]
4. Default steady state (backlogs $\approx 0$, calm demand) & all $.375$ & Cruiser: balanced nominal operating point \\
\bottomrule
\end{tabular}
\end{table}

To enforce flow conservation under severe starvation, an order floor guardrail $a_t^i = \min(20, \max(a_t^i, \text{incoming}))$ is evaluated across all controllers whenever a tier's inventory is zero and backlog is positive, ensuring fully starved echelons order at least the incoming demand.

%=======================================================================
\section{Experimental Setup}

\begin{table}[t]
\caption{Environment and training configuration}
\label{tab:config}
\centering
\scriptsize
\setlength{\tabcolsep}{4pt}
\begin{tabular}{@{}llll@{}}
\toprule
Parameter & Value & Parameter & Value \\
\midrule
Echelons $M$ & 4 & Algorithm & HAPPO \\
Lead time $L$ & 4 & Actor & FC--GRU(128,2)--FC \\
Horizon $T$ & 200 & Critic input & 28-D global \\
$I_0, P_0, B_0$ & 10, 10, 0 & Optimizer & Adam, $10^{-4}$ \\
Action space & $\{0,\dots,20\}$ & $\gamma$, GAE $\lambda$ & 0.95, 0.95 \\
$h$, $b$ & 1.0, 1.0 & PPO clip $\epsilon$ & 0.20 \\
Actor obs. & 7-D / 8-D & Entropy $\beta$ & 0.01 \\
Intervention $N$ & 20 & Value loss & Huber, $\delta{=}10$ \\
Curriculum & 70/15/15 & Rollout threads & 5 \\
Hardware & RTX A5000 & Env.\ steps & $3 \times 10^6$ \\
\bottomrule
\end{tabular}
\end{table}

Table~\ref{tab:config} lists the configuration. Actors and critic are
recurrent~\cite{chung2014} so policies condition on demand history. Both
policies train for $3\times10^6$ environment steps and are evaluated
deterministically via $\arg\max$. The three supervisors of
Section~\ref{sec:res-sup} drive the \emph{same} frozen weights on identical
traces.

\textbf{Reproduction fidelity.} Our baseline is not the numbers printed
in~\cite{liu2022} but the authors' released
implementation~\cite{liu2022code}, trained and evaluated by us under
Table~\ref{tab:config}, and the two do not coincide. The paper reports a
per-period cost of $34.00$, average fluctuation $0.758$ and bullwhip
amplification $92.37\%$ for the four-echelon chain; the same code in our
environment gives $83.90$, $0.863$ and $82.2\%$. We cannot account for the
gap and do not claim to have reproduced the published figures. Our claim is
narrower: every policy compared here shares one codebase, environment,
configuration and trace set, so the comparison between them is internally
valid. The repository ships its serial environment with an echelon count of
three and $\alpha = 0.5$; we set these to four and $0.375$, the
value~\cite{liu2022} reports for that topology, with the common stockout floor
applied during evaluation.

\textbf{Suites.} Two hundred and eighty frozen traces of $T=200$, $56{,}000$
periods in total, generated once from fixed seeds and reused by every
controller. \textbf{S1} ($100$ traces) is Merton jump-diffusion
demand~\cite{merton1976} simulated as in~\cite{glasserman2004}, matching the
training distribution, with empirical mean $8.97$ and standard deviation
$5.06$ (while training clamps demand to at least $1$, S1 evaluation traces
contain $23$ zero-demand periods; both generators cap demand at $20$).
Disruption suites S2--S4 are generated from clipped Poisson draws:
\textbf{S2} ($30$) holds a calm baseline of mean
$10.02$ until $t=60$, applies a $15$-period pulse drawn uniformly on
$[35,45]$ with mean $39.99$, four times baseline and twice the $20$-unit
order cap, then returns to $9.98$ from $t=75$. \textbf{S3} ($120$) steps
demand permanently upward at $t=80$, from Poisson mean $10$ ($9.95$ empirical)
to $19$ ($18.94$ empirical, $+90.5\%$), settling just under the order cap.
\textbf{S4} ($30$) collapses it permanently at $t=80$, from a high-load Poisson
mean of $18$ ($17.42$ empirical) to $5$ ($4.93$ empirical, $-71.7\%$).
All figures are empirical means over the frozen traces,
not generator targets. S2 and S3 open inside the training distribution and
leave it when the shock fires; S4 is outside it throughout, its pre-shock
mean exceeding every S1 trace mean before collapsing below the S1 mean. S3
gets the largest allocation because it is where the controllers were closest
at smaller sample sizes. Whether demand sits above the $20$-unit order cap,
at it, or well below it is what separates the three disruption regimes.

\textbf{Metrics.} Per-period cost is
$\bar C = \tfrac{1}{T}\sum_t \sum_i (I_t^i + B_t^i)$. Order fluctuation is
the coefficient of variation $F_i = \sigma(a^i)/\mu(a^i)$
following~\cite{fransoo2000}, with $\bar F$ its mean and bullwhip
amplification $A = (F_3/F_0 - 1)$. Fill rate is the share of effective
demand, $D_t + B_{t-1}^0$, met from stock on hand. Service level is evaluated
as $1 - (\sum_t B_t^0 / \sum_t D_t)$, measuring demand fulfilled without
carried backlog. Recovery time $T_{\text{rec}}$ measures the number of terminal
backlog-free periods ($\sum_i B_t^i < 0.5$) sustained at the end of the episode;
under sustained step-up (S3) capacity saturation drives $T_{\text{rec}}$ to zero for
all controllers. Comparisons are paired $t$-tests on the shared traces.

%=======================================================================
\section{Results}

\subsection{What the curriculum buys}
\label{sec:res-curr}

\begin{table}[t]
\caption{Headline metrics over 280 frozen traces. ``Base'' is our
reproduction of the released implementation at fixed $\alpha=0.375$; the
remaining three columns drive identical tri-state weights under different
supervisors. Best per row in bold. Service level and recovery time are
reported for S1 in the text.}
\label{tab:main}
\centering
\scriptsize
\setlength{\tabcolsep}{4pt}
\begin{tabular}{@{}llcccc@{}}
\toprule
Suite & Metric & Base & Static & Rules & LLM \\
\midrule
\multirow{3}{*}{\shortstack[l]{S1\\in-dist.\\(100)}}
 & Cost              & 83.90 & \textbf{82.78} & 85.23 & 84.86 \\
 & Fill rate (\%)    & 55.7 & 62.8 & \textbf{63.9} & 63.1 \\
 & Bullwhip $A$ (\%) & \textbf{+82.2} & +91.1 & +84.1 & +89.3 \\
\midrule
\multirow{3}{*}{\shortstack[l]{S2\\pulse\\(30)}}
 & Cost              & 181.71 & \textbf{158.37} & 164.24 & 161.35 \\
 & Fill rate (\%)    & 10.0 & \textbf{13.4} & 13.3 & 13.3 \\
 & Bullwhip $A$ (\%) & +93.5 & +79.5 & \textbf{+74.6} & +81.7 \\
\midrule
\multirow{3}{*}{\shortstack[l]{S3\\step-up\\(120)}}
 & Cost              & \textbf{141.10} & 150.63 & 155.65 & 154.19 \\
 & Fill rate (\%)    & \textbf{21.1} & 19.8 & 19.7 & 19.5 \\
 & Bullwhip $A$ (\%) & \textbf{+37.2} & +45.2 & +42.9 & +45.5 \\
\midrule
\multirow{3}{*}{\shortstack[l]{S4\\step-down\\(30)}}
 & Cost              & 135.96 & 126.57 & \textbf{121.81} & 123.26 \\
 & Fill rate (\%)    & 26.4 & 28.4 & \textbf{28.8} & \textbf{28.8} \\
 & Bullwhip $A$ (\%) & \textbf{+26.4} & +34.0 & +32.4 & +33.8 \\
\bottomrule
\end{tabular}
\end{table}

Table~\ref{tab:main} reports the four suites and Table~\ref{tab:stats} the
paired differences. The \textsc{Static} column isolates the curriculum: it
holds $\alpha = 0.375$ at every tier for the whole episode, the same nominal
operating point as the reproduced specialist, so any difference between the
two is attributable to how the weights were trained.

Under the pulse shock cost falls by $23.35$ units, $12.8\%$; under the
demand collapse by $9.38$, $6.9\%$. Both survive Holm correction across the
sixteen cost tests reported here and are the paper's two robust cost
effects. In distribution the two policies are indistinguishable on cost
($p = 0.45$) but not on availability: fill rate rises $7.10$ points, from
$55.7\%$ to $62.8\%$, with interval $[5.37, 8.84]$, and service level from
$31.0\%$ to $43.3\%$. Seven more orders in a hundred are met from stock at
no cost premium. Fill rate also rises under the pulse and the collapse, both
at $p < 10^{-6}$. This, not cost, is the curriculum's most strongly
supported effect.

Bullwhip amplification moves the other way and we report it plainly: the
curriculum reduces it under the pulse, $+93.5\%$ to $+79.5\%$, but raises it
by seven to nine points in the other three suites. The policy orders more
actively to hold availability up, and with $h=b$ that surfaces as
amplification rather than cost, a trade whose worth depends on a cost
structure this benchmark does not express.

\subsection{Where the curriculum fails}

Under the sustained step-up the specialist wins on both axes: cost $+9.53$
or $6.8\%$ in its favour ($p = 0.015$) and fill rate $21.1\%$ against
$19.8\%$ ($p = 0.009$). We treat this as a boundary condition rather than
noise.

The mechanism is visible in the inventory. Post-shock mean demand in S3 is
$18.94$ against a $20$-unit order cap, so the chain is capacity-bound for
much of the horizon: downstream backlog accumulates under any ordering
policy. Cumulative post-shock holding cost is $290.8$ for the specialist and
$515.9$ for the tri-state policy. The curriculum teaches the policy to carry
buffer: exactly what pays under a transient shock or a collapse, and
exactly what becomes dead weight when the bottleneck is structural. The
specialist is cheaper because it does not try.

Two caveats bound this. The cost difference does not survive Holm correction
over the sixteen tests, so its magnitude is provisional, though the
direction is corroborated by the holding-cost split and the fill-rate
result. And the effect requires a cap that binds.

\subsection{Ablating the supervisor}
\label{sec:res-sup}

\begin{table}[t]
\caption{Paired differences on the shared traces, $\Delta$ with $p$ in
parentheses. Negative favours the first-named controller. Only the two
entries in the first row marked with $^{\ast}$ survive Holm correction
across the sixteen cost tests.}
\label{tab:stats}
\centering
\scriptsize
\setlength{\tabcolsep}{3pt}
\begin{tabular}{@{}lcccc@{}}
\toprule
Contrast & S1 (100) & S2 (30) & S3 (120) & S4 (30) \\
\midrule
\multicolumn{5}{@{}l}{\emph{Per-period cost}}\\
Static $-$ Base  & $-1.13$ & $-23.35^{\ast}$ & $+9.53$ & $-9.38^{\ast}$ \\
                 & (.45) & ($4{\times}10^{-7}$) & (.015) & ($2{\times}10^{-4}$) \\
\addlinespace[2pt]
Rules $-$ Static & $+2.45$ (.044) & $+5.88$ (.062) & $+5.02$ (.16) & $-4.77$ (.005) \\
LLM $-$ Static   & $+2.08$ (.14)  & $+2.99$ (.31)  & $+3.57$ (.32) & $-3.31$ (.054) \\
LLM $-$ Rules    & $-0.37$ (.79)  & $-2.89$ (.39)  & $-1.45$ (.70) & $+1.45$ (.31) \\
\midrule
\multicolumn{5}{@{}l}{\emph{Customer fill rate, percentage points}}\\
Static $-$ Base  & $+7.10$ & $+3.40$ & $-1.26$ & $+2.05$ \\
                 & ($1{\times}10^{-12}$) & ($1{\times}10^{-11}$) & (.009) & ($2{\times}10^{-7}$) \\
\bottomrule
\end{tabular}
\end{table}

\begin{figure}[t]
\centering
\includegraphics[width=\columnwidth]{figures/effects.pdf}
\caption{Paired effects with $95\%$ confidence intervals. (a) The
curriculum, isolated by holding $\alpha$ fixed: it pays under the pulse and
the collapse and costs under the step-up. (b) The twelve supervisor
contrasts on identical frozen weights; only the S4 rules-over-static
interval excludes zero on the helpful side, and the S1 rules interval
excludes it on the harmful side.}
\label{fig:effects}
\end{figure}

The three supervisors drive identical weights, so any difference between
them is the supervisor's own contribution; Fig.~\ref{fig:effects}(b) gives
all twelve paired contrasts.

The language model and the rule base are indistinguishable everywhere,
$p \geq 0.31$ in all four suites. On S1, where $n=100$, the paired
difference is $-0.37$ with interval $[-3.07, +2.33]$. That is not merely a
failure to reject: at this sample size it bounds any true difference between
a frontier language model and a hand-tuned rule base to roughly $\pm 3$ cost
units against a mean near $85$. We read it as an equivalence result.

Run-time switching against a static $\alpha$ produces two nominal effects
with opposite signs. Under the collapse the rule base is cheaper by $4.77$,
interval $[-7.96, -1.57]$, and the language model captures most of that
gain; in distribution the rule base is \emph{more} expensive by $2.45$.
Neither survives correction over the full family, but together they suggest
that adaptivity earns its keep when there is a regime change to respond to
and costs a little when there is not. Terminal stability further corroborates the supervisory mechanism:
on S1, the adaptive rule base sustains an average of $6.66$ backlog-free
periods at the close of the episode compared to $4.52$ for the static arm,
demonstrating that dynamic intervention clears transient disturbances and
restores steady state earlier. Under the step-up (S3), continuous capacity
saturation prevents any controller from reaching a clean terminal state ($T_{\text{rec}} = 0$).

The language model's contribution is therefore neither speed nor cost, but
that it reaches rule-based quality without manual policy threshold engineering,
guided by structured telemetry indicators, and that each decision arrives with
an auditable diagnosis. In one logged
window after the pulse subsided, demand had normalized but upstream
amplification remained high; the model diagnosed ``an upstream overflow
problem'' and held \textsc{Shock Absorber} to let pipeline stock drain
rather than re-ordering. The rule base reaches the same action, but only
because a threshold was tuned to this environment in advance, and it emits
no account of why. Because the two are equivalent, the headline claims
reproduce from the rule-based arm alone. This matters: \texttt{gpt-5.6-luna}
was a preview endpoint that may not remain generally available, and the rule
base is a local, dependency-free substitute indistinguishable from it.

%=======================================================================
\section{Discussion and Limitations}

Two results surprised us. The first is where the robustness comes from: we
built the architecture around run-time regime switching, and the ablation
shows the curriculum already carries most of the gain, with switching adding
to it mainly under contraction. A deployment that cannot afford a
supervisory call in the loop can therefore take the curriculum's benefit for
free, and one that can afford it knows where the extra lies. It is also why
we report the fixed arm: an LLM-in-the-loop system evaluated without one
cannot say which component earned the result. The second surprise is the
regime in which the curriculum loses, which gives the mechanism a
falsifiable shape: buffer is an asset when the constraint is transient and
a liability when it is structural.

Several limitations bound the result.
\textbf{Cost symmetry.} All results assume $h = b = 1$, the benchmark's
structure, adopted so every policy faces identical economics. Since the
central trade is inventory against backlog, the conclusions are conditional
on that ratio, and the buffer mechanism behind both the S2 gain and the S3
loss makes a $b/h$ sweep the experiment we would most want.
\textbf{Action saturation.} Orders cap at $20$ units per echelon per period.
Under the pulse, demand averages $39.99$ for fifteen periods, twice what any
tier can order; under the step-up it settles at $18.94$, just below the cap.
A large share of disruption demand is therefore structurally unfillable,
which is why absolute fill rates are low, and the binding cap is a
precondition for the S3 reversal. These figures measure behaviour against a
capacity wall, not behaviour with slack.
\textbf{Statistical resolution and multiplicity.} At $n=100$ and $n=120$ the
supervisor contrasts on S1 and S3 support an equivalence reading; at $n=30$,
S2 and S4 do not, and an effect the size of S4's could be missed elsewhere.
Of sixteen cost tests only two survive Holm correction, so the remaining
nominal effects are descriptive rather than confirmatory.
\textbf{Latency and dependency.} Each supervisory call adds
$0.5$--$2.5$\,s and, for the LLM arm, an external API; this suits planning
horizons of hours or days, not sub-second control.
\textbf{Topology.} The telemetry schema and prompt target a four-echelon
serial chain; a general network would need graph-structured state or a
hierarchy of supervisors.

The next step follows from the ablation. If adaptivity rather than its
implementation is what matters, the question is whether an LLM supervisor
transfers to a topology whose thresholds have \emph{not} been hand-tuned,
where the rule base would have to be rewritten and the model would not. We
have not tested that.

%=======================================================================
\section{Conclusion}

We revisited the cooperation weight in MADRL-based multi-echelon inventory
control and made it a control input rather than a tuned constant. Training
one HAPPO policy across three cooperation regimes cuts per-period cost by
$12.8\%$ under a pulse shock and $6.9\%$ under a demand collapse relative to
a reproduced single-$\alpha$ specialist, and raises in-distribution fill
rate by $7.1$ points at cost parity. All of this holds with $\alpha$ frozen
at execution time, so the gain needs no run-time machinery to collect. The robustness is
uneven: under a sustained step-up that drives demand to the order cap the
curriculum costs $6.8\%$ more, because the buffer that pays under a
transient shock is dead weight when the bottleneck is structural. Ablating
the supervisor on identical weights, a language model given only the
topology and the regimes matches a hand-tuned rule base to within $\pm 3$
cost units at $n=100$, reaching expert heuristic quality with no threshold
engineering; switching pays most under contraction. Reporting a fixed-$\alpha$
arm is what makes these contributions separable, and we would encourage its
inclusion in comparable studies.

%=======================================================================
\begin{thebibliography}{00}

\bibitem{lee1997} H. L. Lee, V. Padmanabhan, and S. Whang, ``The bullwhip effect in supply chains,'' \emph{Sloan Manage. Rev.}, vol. 38, no. 3, pp. 93--102, 1997.

\bibitem{clark1960} A. J. Clark and H. Scarf, ``Optimal policies for a multi-echelon inventory problem,'' \emph{Manage. Sci.}, vol. 6, no. 4, pp. 475--490, 1960.

\bibitem{shang2003} K. H. Shang and J.-S. Song, ``Newsvendor bounds and heuristic for optimal policies in serial supply chains,'' \emph{Manage. Sci.}, vol. 49, no. 5, pp. 618--638, 2003.

\bibitem{boute2022} R. N. Boute, J. Gijsbrechts, W. van Jaarsveld, and N. Vanvuchelen, ``Deep reinforcement learning for inventory control: A roadmap,'' \emph{Eur. J. Oper. Res.}, vol. 298, no. 2, pp. 401--412, 2022.

\bibitem{gijsbrechts2021} J. Gijsbrechts, R. N. Boute, J. A. Van Mieghem, and D. J. Zhang, ``Can deep reinforcement learning improve inventory management? Performance on lost sales, dual-sourcing, and multi-echelon problems,'' \emph{Manuf. Serv. Oper. Manage.}, vol. 24, no. 3, pp. 1349--1368, 2022.

\bibitem{harsha2022} P. Harsha, A. Jagmohan, J. Kalagnanam, B. Quanz, and D. Singhvi, ``Deep policy iteration with integer programming for inventory management,'' \emph{Manuf. Serv. Oper. Manage.}, vol. 27, no. 2, pp. 369--388, 2025.

\bibitem{vosooghidizaji2020} M. Vosooghidizaji, A. Taghipour, and B. Canel-Depitre, ``Supply chain coordination under information asymmetry: A review,'' \emph{Int. J. Prod. Res.}, vol. 58, no. 6, pp. 1805--1834, 2020.

\bibitem{liu2022} X. Liu, M. Hu, Y. Peng, and Y. Yang, ``Multi-agent deep reinforcement learning for multi-echelon inventory management,'' \emph{Prod. Oper. Manage.}, vol. 34, no. 7, pp. 1836--1856, 2025, doi: 10.1177/10591478241305863.

\bibitem{liu2022code} X. Liu, ``Multi-agent deep reinforcement learning on multi-echelon inventory management,'' GitHub repository, commit a7e5a3e. [Online]. Available: \url{https://github.com/xiaotianliu01/Multi-Agent-Deep-Reinforcement-Learning-on-Multi-Echelon-Inventory-Management} (accessed Sep. 29, 2026).

\bibitem{kuba2021} J. G. Kuba, R. Chen, M. Wen, Y. Wen, F. Sun, J. Wang, and Y. Yang, ``Trust region policy optimisation in multi-agent reinforcement learning,'' in \emph{Proc. Int. Conf. Learn. Represent. (ICLR)}, 2022.

\bibitem{vanvuchelen2020} N. Vanvuchelen, J. Gijsbrechts, and R. Boute, ``Use of proximal policy optimization for the joint replenishment problem,'' \emph{Comput. Ind.}, vol. 119, 103239, 2020.

\bibitem{schulman2017} J. Schulman, F. Wolski, P. Dhariwal, A. Radford, and O. Klimov, ``Proximal policy optimization algorithms,'' arXiv:1707.06347, 2017.

\bibitem{oroojlooy2022} A. Oroojlooyjadid, M. Nazari, L. V. Snyder, and M. Tak\'{a}\v{c}, ``A deep Q-network for the beer game: Deep reinforcement learning for inventory optimization,'' \emph{Manuf. Serv. Oper. Manage.}, vol. 24, no. 1, pp. 285--304, 2022.

\bibitem{lowe2017} R. Lowe, Y. Wu, A. Tamar, J. Harb, P. Abbeel, and I. Mordatch, ``Multi-agent actor-critic for mixed cooperative-competitive environments,'' in \emph{Adv. Neural Inf. Process. Syst. (NeurIPS)}, vol. 30, 2017.

\bibitem{yu2021} C. Yu, A. Velu, E. Vinitsky, J. Gao, Y. Wang, A. Bayen, and Y. Wu, ``The surprising effectiveness of PPO in cooperative, multi-agent games,'' in \emph{Adv. Neural Inf. Process. Syst. (NeurIPS) Datasets and Benchmarks Track}, 2022.

\bibitem{ahn2022} M. Ahn \emph{et al.}, ``Do as I can, not as I say: Grounding language in robotic affordances,'' arXiv:2204.01691, 2022.

\bibitem{yao2023} S. Yao, J. Zhao, D. Yu, N. Du, I. Shafran, K. Narasimhan, and Y. Cao, ``ReAct: Synergizing reasoning and acting in language models,'' in \emph{Proc. Int. Conf. Learn. Represent. (ICLR)}, 2023.

\bibitem{kwon2023} M. Kwon, S. M. Xie, K. Bullard, and D. Sadigh, ``Reward design with language models,'' in \emph{Proc. Int. Conf. Learn. Represent. (ICLR)}, 2023.

\bibitem{ma2024} Y. J. Ma \emph{et al.}, ``Eureka: Human-level reward design via coding large language models,'' in \emph{Proc. Int. Conf. Learn. Represent. (ICLR)}, 2024.

\bibitem{littman1994} M. L. Littman, ``Markov games as a framework for multi-agent reinforcement learning,'' in \emph{Proc. Int. Conf. Mach. Learn. (ICML)}, 1994, pp. 157--163.

\bibitem{merton1976} R. C. Merton, ``Option pricing when underlying stock returns are discontinuous,'' \emph{J. Financ. Econ.}, vol. 3, no. 1--2, pp. 125--144, 1976.

\bibitem{glasserman2004} P. Glasserman, \emph{Monte Carlo Methods in Financial Engineering}. New York, NY, USA: Springer, 2004.

\bibitem{fransoo2000} J. C. Fransoo and M. J. F. Wouters, ``Measuring the bullwhip effect in the supply chain,'' \emph{Supply Chain Manage.}, vol. 2, no. 2, pp. 78--89, 2000.

\bibitem{sterman1989} J. D. Sterman, ``Modeling managerial behavior: Misperceptions of feedback in a dynamic decision making experiment,'' \emph{Manage. Sci.}, vol. 35, no. 3, pp. 321--339, 1989.

\bibitem{zipkin2000} P. H. Zipkin, \emph{Foundations of Inventory Management}. New York, NY, USA: McGraw-Hill, 2000.

\bibitem{bengio2009} Y. Bengio, J. Louradour, R. Collobert, and J. Weston, ``Curriculum learning,'' in \emph{Proc. Int. Conf. Mach. Learn. (ICML)}, 2009, pp. 41--48.

\bibitem{tobin2017} J. Tobin, R. Fong, A. Ray, J. Schneider, W. Zaremba, and P. Abbeel, ``Domain randomization for transferring deep neural networks from simulation to the real world,'' in \emph{Proc. IEEE/RSJ Int. Conf. Intell. Robots Syst. (IROS)}, 2017, pp. 23--30.

\bibitem{kirk2023} R. Kirk, A. Zhang, E. Grefenstette, and T. Rockt\"{a}schel, ``A survey of zero-shot generalisation in deep reinforcement learning,'' \emph{J. Artif. Intell. Res.}, vol. 76, pp. 201--264, 2023.

\bibitem{yang2024} C. Yang, X. Wang, Y. Lu, H. Liu, Q. V. Le, D. Zhou, and X. Chen, ``Large language models as optimizers,'' in \emph{Proc. Int. Conf. Learn. Represent. (ICLR)}, 2024.

\bibitem{chung2014} J. Chung, C. Gulcehre, K. Cho, and Y. Bengio, ``Empirical evaluation of gated recurrent neural networks on sequence modeling,'' arXiv:1412.3555, 2014.

\end{thebibliography}

\end{document}
